\section{Summary and recommendations}
\label{sec:summary}

\subsection{What was established}

\begin{enumerate}[leftmargin=1.6em,itemsep=4pt]
  \item \textbf{Models collapse onto spectra.} Classified by \emph{where the bump appears} rather
than by theory, and with lepton flavour and $b$-tag content counted as part of the mass axis, the
public BSM resonance model classes map onto \textbf{46 canonical spectra}. The number of searches is
set by the number of spectra, not by the number of models.
  \item \textbf{The search budget.} $\Nsig \approx 3.7\times10^{3}$ inclusively,
        $6.6\times10^{3}$ at realistic event-selection granularity, hence $\Zloc \approx 6.4$ to
$6.5$ for a $5\sigma$ global discovery, stable to $\pm0.1\sigma$ against every counting choice
varied. \textbf{Covering the entire known BSM resonance space costs about one extra unit of $\sigma$
over a single counting experiment.}
  \item \textbf{The combinatorial scan costs half a sigma more.} 150 exclusive object categories
        and 1094 mass groups give $\Nsig = 1.4\times10^{5}$ and $\Zloc = 6.98$.
  \item \textbf{Most of the reachable space has never been scanned.} 57 of the 82 reachable object
compositions have no published ATLAS bump hunt, under a deliberately generous mapping. Two-body
space is saturated, while three- and four-body space is 75 to 90\,\% untouched. The ATLAS general
search declined that coverage explicitly on look-elsewhere grounds, a trade the numbers above
reprice.
  \item \textbf{The arithmetic is self-consistent.} The same trials-factor calculation post-dicts
        the observed population of ATLAS excesses: tens of $3\sigma$, all faded, and no spurious
        $5\sigma$.
  \item \textbf{Rank-based candidate selection is dominated.} At an identical false-alarm budget, a
        fixed threshold confirms more signal than the argmax at every signal strength, and a
        Benjamini--Hochberg rule lands between them. With an imperfect significance estimator the
        argument gets stronger, not weaker: the threshold's margin grows to more than ten
        percentage points and the argmax loses control of its false-alarm rate entirely.
  \item \textbf{A/B splitting relocates the look-elsewhere effect rather than reducing it.} The cost
factorises exactly, and the residual two-coin penalty is $+0.3$ to $+0.5\sigma$ at any trials
count. Symmetrising the scheme --- running both directions and claiming on the union --- does not
recover it: it repairs a $50/50$ split but is a wash against the asymmetric optimum, and it
forfeits procedural blindness. The scheme is worth adopting when the trials factor one would have
to defend exceeds roughly $13\times$ the true one, which is the situation a machine-learning-driven
scan creates.
\end{enumerate}

\subsection{Recommendations}

\paragraph{For an analysis.} Quote the global significance against a counted trials factor and
state the counting convention (published scan windows, resolution elements) explicitly. That is
what makes $6.4$ rather than $5$ the number to beat, and it is defensible in a way that ``we
scanned a lot of spectra'' is not. Select stage-1 candidates on absolute significance at a fixed
expected-fake budget rather than by rank, and quote the \emph{achieved} false-alarm rate rather than
a nominal FDR parameter. If a two-stage scheme is adopted, use an asymmetric split
($f \approx 0.25$ to $0.35$), pre-register the windows and the outcome ladder, use cross-half
trainings, and expect to pay $0.5\sigma$ for the auditability.

\paragraph{For the search program.} The three- and four-body mixed-flavour space is the largest
identified gap, and it is not expensive to close: moving from the published program to a full
combinatorial scan costs $0.5\sigma$ in the discovery bar. Signatures scanned once at Run-1
luminosity and then abandoned, such as $m(Zb)$, cost nothing in trials to revisit, because
Section~\ref{sec:budget} has already counted their axes.

\paragraph{For a sample program.} The exchange rate is unambiguous, whoever is doing the
prioritising. Adding a model to an already-scanned spectrum costs \emph{zero} trials, while opening
a new spectrum costs its $n_s$ --- from $14$ looks for $m(Ht)$ to $411$ for $m(\gamma\gamma)$
(Table~\ref{tab:budget}). Any argument about which channels to invest in should be conducted in
those units.

\subsection{Known limitations}

The fractional resolutions are coarse per-channel central values, which is why every headline
carries a factor-two band. Scan windows are hand-curated from published search families and
inevitably involve judgement at the edges. The trials factor for the full ATLAS program used in
Section~\ref{sec:excess} spans a factor ten and is the dominant uncertainty there, although the
conclusion survives the whole range. Cross-channel correlations are neglected, which makes the sum
mildly conservative, and event selections within one search family are partly correlated and often
statistically combined, so the selections-level count is an over-count rather than a refinement.
The model-to-spectrum map is assembled from public model databases and the published search record;
it is complete to the best of our knowledge but is not machine-verifiable, and a missing model class
that populates only an already-counted spectrum would not change $\Nsig$ at all.
